\section{Proofs for Section~\ref{sec:regions}}
\label{app:regions}

The notation is that of the main text, and symbols are recalled where they reappear far from their definition. We
write $B_\Hh$ for the closed unit ball of $\Hh$ and $\mathcal G_\alpha(P)=\{g:0\le g\le\beta,\ \E_P[g]=1\}$ for the set
of densities of the $\alpha$-trimmings of $P$ with respect to $P$, so that
$\mathcal C_\alpha(P)=\{\E_P[g\phi]:g\in\mathcal G_\alpha(P)\}$, and for independent Rademacher signs $\varepsilon_i$
we put $e_k=\E|\varepsilon_1+\dots+\varepsilon_k|/k\le k^{-1/2}$. Several proofs use the separation
$\mathrm{sep}(f;A,B)=\inf_{a\in A}\langle f,a\rangle_\Hh-\sup_{b\in B}\langle f,b\rangle_\Hh$ of
Section~\ref{sec:ramp} through the elementary fact, proved in Lemma~\ref{lem:separation}, that
$\dist(A,B)=\max_{f\in B_\Hh}\mathrm{sep}(f;A,B)$ for convex weakly compact $A$ and $B$, which is \eqref{eq:distmax}.
Proposition~\ref{prop:bangbang} and the uniform bound \eqref{eq:hausdorff} of Section~\ref{sec:regions} are proved
later, in Section~\ref{app:trimming}, next to the tools they need.

Two technical points are settled once and for all. By (A1) and (A2) of Assumption~\ref{ass:regularity}, $\phi$ is
weakly measurable with $\|\phi\|_\Hh\le\kappa$, hence Bochner integrable by Pettis' theorem, so $\E_P[g\phi]$ is well
defined for bounded measurable $g$. And the suprema over directions and cuts that appear below are measurable. Indeed,
$|f(x)-f'(x)|\le\kappa\|f-f'\|_\Hh$ and $|(u-t)_+-(u-t')_+|\le|t-t'|$ make $(f,t)\mapsto(f(x)-t)_+$ jointly
continuous, and dominated convergence makes $f\mapsto\E_P[(f-t)_+]$ continuous, so every such supremum may be taken over
a countable dense subset without changing its value.

\subsection{The results of Section~\ref{sec:regions}}

We first prove \eqref{eq:tvintersect}. A law $R$ is a common $\alpha$-trimming exactly when $(1-\alpha)R$ lies below
both $P$ and $Q$, so the question is one of mass. Let $P\wedge Q$ be the largest measure below both $P$ and $Q$; it has
density $\min(p,q)$ with respect to any measure dominating both, and mass $1-\TV(P,Q)$. If $\TV(P,Q)\le\alpha<1$, the
law $R=(P\wedge Q)/(1-\TV(P,Q))$ has densities at most $1/(1-\TV(P,Q))\le\beta$ with respect to $P$ and to $Q$, so it
lies in $R_\alpha(P)\cap R_\alpha(Q)$. Conversely, if $R$ lies in both sets, then $(1-\alpha)R\le P$ and
$(1-\alpha)R\le Q$, hence $(1-\alpha)R\le P\wedge Q$ by maximality of $P\wedge Q$, and comparing the masses of the two
sides gives $1-\alpha\le1-\TV(P,Q)$.

The same comparison of masses shows that laws which lose different fractions $\alpha_1$ and $\alpha_2$ and agree on what
remains must have $\alpha_1\vee\alpha_2\ge\alpha^\star$. Their common remainder is a probability measure with density
at most $\min\{p/(1-\alpha_1),q/(1-\alpha_2)\}\le\min(p,q)/(1-\alpha_1\vee\alpha_2)$, whose integral is
$(1-\alpha^\star)/(1-\alpha_1\vee\alpha_2)$; as it must be at least one, $\alpha_1\vee\alpha_2\ge\alpha^\star$.

\begin{proof}[Proof of Proposition~\ref{prop:georep}]
The map $g\mapsto g\cdot P$ is a bijection from $\mathcal G_\alpha(P)$ onto $R_\alpha(P)$, and Bochner
integrability gives $\mu_{g\cdot P}=\E_P[g\phi]$, so $P'\mapsto\mu_{P'}$ maps $R_\alpha(P)$ onto
$\mathcal C_\alpha(P)$. The infimum in \eqref{eq:trimMMDpop} is therefore the infimum of $\|u-v\|_\Hh$ over
$u\in\mathcal C_\alpha(P)$ and $v\in\mathcal C_\alpha(Q)$, which is \eqref{eq:setdist}.

For $P=P_n$ the only care needed is with sample points that no measurable set separates, so we group the sample points
into such classes. A density with respect to $P_n$ is measurable, hence constant on each class, and so defines a capped
weight vector $w_i=g(X_i)/n$ with the same embedding. Conversely, averaging a capped weight vector within each class
preserves the cap, the normalization and $\sum_iw_i\phi(X_i)$, because $k(\cdot,x')$ is measurable for every $x'$ by
(A1), so that $\phi$ is constant on a class; the averaged vector is a density. The set is a polytope, as the
image of the capped simplex, itself a polytope, under the linear map $w\mapsto\sum_iw_i\phi(X_i)$.
\end{proof}

\begin{proof}[Proof of Lemma~\ref{lem:attainment}]
Everything except norm-compactness follows from the weak compactness of the regions.

Convexity and the bound are immediate: $\mathcal G_\alpha(P)$ is convex and $g\mapsto\E_P[g\phi]$ is linear, and
$\|\E_P[g\phi]\|_\Hh\le\E_P[g\|\phi\|_\Hh]\le\kappa$. For compactness, regard $\mathcal G_\alpha(P)$ as a subset of
$L^\infty(P)=(L^1(P))^*$. It is bounded, and it is weak$^*$-closed, since each constraint has the form
$\int gh\,dP\ \{\ge,\le,=\}\ c$ for a fixed $h\in L^1(P)$, with $h=\mathbf 1_A$ for the pointwise bounds and $h\equiv1$
for the normalization; so it is weak$^*$-compact by the Banach--Alaoglu theorem. The map $T:g\mapsto\E_P[g\phi]$ is
weak$^*$-to-weak continuous, because $\langle Tg,e\rangle_\Hh=\int ge\,dP$ for $e\in\Hh$, and $e$ is a bounded
function, $|e(x)|\le\kappa\|e\|_\Hh$. So $\mathcal C_\alpha(P)$ is weakly compact, hence weakly
closed and norm-closed. It is even norm-compact, because the embedding is a Hilbert--Schmidt operator. The adjoint of
$T:L^2(P)\to\Hh$, $Tg=\E_P[g\phi]$, sends $h\in\Hh$ to the function $h$ in $L^2(P)$. For an orthonormal basis $(h_j)$
of $\Hh$, Parseval's identity gives $\sum_j\|T^*h_j\|^2_{L^2(P)}=\E_P\sum_jh_j(X)^2=\E_P\|\phi(X)\|_\Hh^2\le\kappa^2$,
so $T^*$, and hence $T$, is Hilbert--Schmidt and therefore compact. The set $\mathcal G_\alpha(P)$ is bounded in
$L^2(P)$, so $\mathcal C_\alpha(P)=T(\mathcal G_\alpha(P))$ is relatively compact, and it is norm-closed.

The norm is weakly lower semicontinuous on the weakly compact product of the two regions, so the distance is attained,
at points $\mu_{P^\star},\mu_{Q^\star}$ for trimmings $P^\star,Q^\star$ by Proposition~\ref{prop:georep}. If
$D_\alpha(P,Q)=0$, the two points coincide, so the intersection is nonempty, and every element of it is
$\mu_{P^\star}=\mu_{Q^\star}$ for some trimmings $P^\star,Q^\star$, which are equal by (A3).
\end{proof}

The equivalence \eqref{eq:zeroeq} follows. If $\TV(P,Q)\le\alpha$, \eqref{eq:tvintersect} gives a common trimming,
whose embedding lies in both regions, so $D_\alpha(P,Q)=0$; this direction uses neither attainment nor (A3).
Conversely, if $D_\alpha(P,Q)=0$, Lemma~\ref{lem:attainment} gives a common trimming, and \eqref{eq:tvintersect} gives
$\TV(P,Q)\le\alpha$.

\begin{proof}[Proof of Lemma~\ref{lem:cvar}]
The first equality holds because $\langle\E_P[g\phi],f\rangle_\Hh=\E_P[gf]$. The second is a bathtub problem, in
which every cut gives an upper bound and the bound at one cut is attained. For any $t$ and $g\in\mathcal G_\alpha(P)$, since
$0\le g\le\beta$ and $\E_Pg=1$,
\[
\E_P[gf]=t+\E_P[g(f-t)]\le t+\beta\,\E_P[(f-t)_+] ,
\]
so the supremum is at most the infimum. For the converse we fill the values of $f$ above a quantile with the largest
weight allowed. Let $t^\star$ be an $\alpha$-quantile of $f(X)$, so that
$\Pp_P(f>t^\star)\le1-\alpha\le\Pp_P(f\ge t^\star)$, and put $g^\star=\beta\mathbf 1\{f>t^\star\}+c\mathbf
1\{f=t^\star\}$ with $c=(1-\beta\Pp_P(f>t^\star))/\Pp_P(f=t^\star)$ when $\Pp_P(f=t^\star)>0$, and $c=0$ otherwise. Multiplied by $\beta$, the quantile inequalities read
$\beta\Pp_P(f>t^\star)\le1\le\beta\Pp_P(f\ge t^\star)$, hence $c\in[0,\beta]$ and $\E_P[g^\star]=1$ (when
$\Pp_P(f=t^\star)=0$ they force $\beta\Pp_P(f>t^\star)=1$). Since $(f-t^\star)_+=(f-t^\star)\mathbf 1\{f>t^\star\}$,
we get $\E_P[g^\star f]=t^\star+\beta\E_P[(f-t^\star)_+]$, so $g^\star$ attains the supremum and $t^\star$ the
infimum. Only boundedness and measurability of $f$ were used.
\end{proof}

\subsection{The cut, continuity in the level and a hinge class}

The cut $t$ in \eqref{eq:cvar}, the threshold of the next lemma, can be confined to $[-\kappa,\kappa]$.

\begin{lemma}[Compact range of the threshold]
\label{lem:tcompact}
Let Assumption~\ref{ass:regularity} hold, $\alpha\in[0,1)$, and $R$ any probability measure. For every $f\in B_\Hh$
the infimum in \eqref{eq:cvar}, with $R$ in place of $P$, is attained at some $t^\star\in[-\kappa,\kappa]$.
\end{lemma}

\begin{proof}
The function $A(t)=t+\beta\E_R[(f(X)-t)_+]$ is convex, with right derivative $1-\beta\Pp_R(f(X)>t)$ by dominated
convergence. Since $|f|\le\kappa$, $A(t)=t$ for $t>\kappa$, while for $t<-\kappa$ the right derivative is
$1-\beta=-\alpha\beta\le0$. A continuous convex function that is nonincreasing on $(-\infty,-\kappa]$ and increasing
on $[\kappa,\infty)$ attains its infimum in $[-\kappa,\kappa]$.
\end{proof}

\begin{lemma}[Local Lipschitz continuity in $\alpha$]
\label{lem:continuity}
Under Assumption~\ref{ass:regularity}, for every $\bar\alpha\in[0,1)$ and $\alpha,\alpha'\in[0,\bar\alpha]$,
\[
\left|D_\alpha(P,Q)-D_{\alpha'}(P,Q)\right|\le\frac{4\kappa}{(1-\bar\alpha)^2}\,|\alpha-\alpha'| .
\]
In particular $\alpha\mapsto D_\alpha(P,Q)$ is continuous on $[0,1)$.
\end{lemma}

\begin{proof}
The level enters the support functions only through $\beta$, as an infimum of functions affine in $\beta$, so it
suffices to bound their slopes. Write $\beta'=(1-\alpha')^{-1}$. For $f\in B_\Hh$, Lemmas~\ref{lem:cvar} and~\ref{lem:tcompact} give
\[
h_{\mathcal C_\alpha(P)}(f)=\inf_{t\in[-\kappa,\kappa]}\big\{t+\beta\,\E_P[(f-t)_+]\big\},
\]
an infimum of affine functions of $\beta$ with slopes $\E_P[(f-t)_+]\in[0,2\kappa]$. An infimum of functions that are
$2\kappa$-Lipschitz is $2\kappa$-Lipschitz, so $|h_{\mathcal C_\alpha(P)}(f)-h_{\mathcal C_{\alpha'}(P)}(f)|\le
2\kappa|\beta-\beta'|$ uniformly in $f$, and likewise for $Q$. By \eqref{eq:distmax} in Section~\ref{app:split},
\[
D_\alpha(P,Q)=\max_{f\in B_\Hh}\big\{-h_{\mathcal C_\alpha(P)}(-f)-h_{\mathcal C_\alpha(Q)}(f)\big\},
\]
and a difference of two maxima is at most the largest difference of the maximized functions, so
$|D_\alpha(P,Q)-D_{\alpha'}(P,Q)|\le4\kappa|\beta-\beta'|$. Finally
$\tfrac{d}{d\alpha}(1-\alpha)^{-1}\le(1-\bar\alpha)^{-2}$ on $[0,\bar\alpha]$, which gives the claim. Only (A1) and (A2)
were used, for arbitrary probability measures, so the lemma applies to $P_n$ and $Q_m$.
\end{proof}

The deviation bounds of Theorem~\ref{thm:twosided}, Corollary~\ref{cor:uniform} and Lemma~\ref{lem:localmod} are
instances of one lemma, in which $n_\star=nm/(n+m)$, as in Section~\ref{sec:analytic}.

\begin{lemma}[A two-sample hinge class]
\label{lem:hinge}
Let (A1) and (A2) of Assumption~\ref{ass:regularity} hold, let $X_1,\dots,X_n$ be i.i.d.\ $P$ and $Y_1,\dots,Y_m$
i.i.d.\ $Q$, the two samples independent, and let $r,\rho_1,\rho_2\ge0$. For each $x\in\mathcal X$ let
$\varphi^1_x,\varphi^2_x:\R\to\R$ be $1$-Lipschitz, with $(x,z)\mapsto\varphi^i_x(z)$ measurable and $\sup_x|\varphi^i_x(0)|<\infty$.
For $(v,a,b)$ in $\mathcal S=\{(v,a,b)\in\Hh\times\R^2:\|v\|_\Hh\le r,\ |a|\le\rho_1,\ |b|\le\rho_2\}$ put
$h^1_{v,a}(x)=\varphi^1_x(a-v(x))$ and $h^2_{v,b}(y)=\varphi^2_y(v(y)-b)$, and let
\[
\Gamma=\sup_{\mathcal S}\big[(P-P_n)h^1_{v,a}+(Q-Q_m)h^2_{v,b}\big],
\qquad
\Gamma^-=\sup_{\mathcal S}\big[(P_n-P)h^1_{v,a}+(Q_m-Q)h^2_{v,b}\big].
\]
Then $\Gamma$ and $\Gamma^-$ are measurable and their means are at most $2[r\kappa/\sqrt{n_\star}+\rho_1e_n+\rho_2e_m]$. If moreover
$|h^1_{v,a}(x)-h^1_{v,a}(x')|\le\Delta_1$ and $|h^2_{v,b}(y)-h^2_{v,b}(y')|\le\Delta_2$ for all points and all
$(v,a,b)\in\mathcal S$, then for every $\chi>0$ each of $\Gamma$ and $\Gamma^-$ exceeds its mean by more than
$\sqrt{\chi(\Delta_1^2/n+\Delta_2^2/m)/2}$ with probability at most $e^{-\chi}$. For one sample the same holds with the
terms of the second sample removed: $n_\star=n$, and no $\rho_2e_m$ and no $\Delta_2^2/m$.
\end{lemma}

\begin{proof}
The mean is bounded by symmetrization, contraction and a bound on the norm of a Rademacher sum in $\Hh$, and the
deviation by McDiarmid's inequality. Measurability is settled first: for fixed samples the bracket is continuous in
$(v,a,b)$, and the population terms are continuous by dominated convergence with the envelope
$\sup_x|\varphi^i_x(0)|+\rho_i+\kappa r$; as $\Hh\times\R^2$ is separable, $\Gamma$ and $\Gamma^-$ are measurable.

For the mean, $\Gamma$ is a supremum of sums of $n+m$ independent centered terms, and symmetrization holds in the signed
form, with no absolute value inside the supremum. We introduce an independent copy of the samples, apply Jensen's
inequality, multiply the differences by independent signs $\varepsilon_i,\varepsilon'_j$, and split the supremum, to get
\[
\E \Gamma\le2\,\E\sup_{\mathcal S}\Big[\frac1n\sum_i\varepsilon_i\varphi^1_{X_i}\big(a-v(X_i)\big)+\frac1m\sum_j\varepsilon'_j
\varphi^2_{Y_j}\big(v(Y_j)-b\big)\Big].
\]
Conditionally on the samples, we remove the maps $\varphi^i$ by the contraction principle for signed Rademacher
averages with a different $1$-Lipschitz map in each coordinate \citep[Lemma~26.9]{shalevshwartzbendavid2014}, applied
to the coordinates $(a-v(X_i))/n$ and $(v(Y_j)-b)/m$ and the maps $z\mapsto\varphi^1_{X_i}(nz)/n$ and
$z\mapsto\varphi^2_{Y_j}(mz)/m$. Its constant is $1$, not the $2$ of the version with absolute values
\citep[Theorem~4.12]{ledouxtalagrand1991}, and it needs no normalization $\varphi(0)=0$. It replaces the bracket by the
linear expression
\[
\langle v,S_\varepsilon\rangle_\Hh+\frac an\sum_i\varepsilon_i-\frac bm\sum_j\varepsilon'_j,
\qquad
S_\varepsilon=\frac1m\sum_j\varepsilon'_j\phi(Y_j)-\frac1n\sum_i\varepsilon_i\phi(X_i),
\]
whose supremum over $\mathcal S$ is $r\|S_\varepsilon\|_\Hh+\rho_1|n^{-1}\sum_i\varepsilon_i|+\rho_2|m^{-1}\sum_j\varepsilon'_j|$.
The means of the two sign averages are $e_n$ and $e_m$ by definition, and
\[
\E\|S_\varepsilon\|_\Hh\le(\E\|S_\varepsilon\|_\Hh^2)^{1/2}\le\kappa/\sqrt{n_\star}
\]
by Jensen's inequality, then independence of the signs and $k(x,x)\le\kappa^2$, since $1/n+1/m=1/n_\star$. This gives the
bound on the mean. The argument for $\Gamma^-$ is the same, with the sample and its copy exchanged.

For the deviation, replacing one $X_i$ moves $\Gamma$ and $\Gamma^-$ by at most $\Delta_1/n$, and one $Y_j$ by at most
$\Delta_2/m$, so McDiarmid's inequality \citep{mcdiarmid1989} gives the tail bound (if $\Delta_1=\Delta_2=0$ every $h$ is
constant and $\Gamma=\Gamma^-=0$).
\end{proof}

\subsection{No informative upper bound}
\label{sec:upperimposs}

The triviality of distribution-free upper bounds on $\TV$ persists on the much
smaller class of pairs whose densities are bounded by $2$, which contains every $\tfrac12$-trimming of the uniform law
and in which no atom can be planted: a procedure valid over the whole class must return the trivial bound, with
probability at least $1-\delta$, even when the two laws are identical.

\begin{proof}[Proof of Theorem~\ref{thm:impossible}]
We hide a pair of laws at total variation one in a fine random partition. Fix an even integer $M\ge2(n+m)$, cut $[0,1]$ into $M$ cells of equal length, and let $S$ be a uniformly random set of
$M/2$ cells. Let $P_S$ and $Q_S$ be the uniform laws on the union of the cells in $S$ and in its complement, with
densities $2\cdot\mathbf 1_S$ and $2\cdot\mathbf 1_{S^c}$. Then $(P_S,Q_S)\in\mathcal P_2$ and $\TV(P_S,Q_S)=1$, so the
coverage requirement gives $\Pp_{P_S,Q_S}(\widehat U\ge1)\ge1-\delta$ for each $S$, and hence under the mixture $\mu_M$
of the laws of the data over the random $S$.

Data drawn from $\mu_M$ are nearly indistinguishable from two uniform samples, because they rarely put two
observations in one cell. Let $E$ be the event that the $n+m$ observations fall in $n+m$ distinct cells. Given $E$, $\mu_M$ and the product
$P_{\mathrm u}^{\otimes(n+m)}$ of uniform laws coincide. Under $\mu_M$, an ordered tuple of distinct cells receives
$X_1,\dots,X_n,Y_1,\dots,Y_m$ with probability $(2/M)^{n+m}$ times the probability that its first $n$ cells lie in $S$
and its last $m$ in $S^c$, which depends only on $n$ and $m$, so under both laws the tuple is uniform given $E$, with
independent uniform positions within the cells. Laws that agree given an event differ in total variation by at most
the sum of the masses of its complement. Under $\mu_M$ an $X$ and a $Y$ never share a cell, and two $X$'s, or two
$Y$'s, do so with probability at most $2/M$; under the product law two observations do so with probability $1/M$. Hence
$\TV(\mu_M,P_{\mathrm u}^{\otimes(n+m)})\le\{n(n-1)+m(m-1)\}/M+(n+m)(n+m-1)/(2M)\le2(n+m)^2/M$, also with an independent
randomizer. Finally,
$\Pp_{P_{\mathrm u},P_{\mathrm u}}(\widehat U\ge1)\ge1-\delta-2(n+m)^2/M$ for every even $M\ge2(n+m)$, and letting
$M\to\infty$ gives the claim.
\end{proof}

Lower bounds escape Theorem~\ref{thm:impossible}, but their rates do not. If $\tilde\alpha$ were a lower confidence
bound valid at level $\delta$ for every pair in $\mathcal P_2$ and also had
$\Pp_{P,Q}(\TV(P,Q)-\tilde\alpha>r)\le\delta'$ for every pair in $\mathcal P_2$, with $r<1$ and $\delta+\delta'<1$,
then $\tilde\alpha+r$ would be an upper bound valid at level $\delta'$, so by Theorem~\ref{thm:impossible} it would
reach $1$ at the uniform pair with probability at least $1-\delta'$, and there $\tilde\alpha\ge1-r>0=\TV$ with
probability larger than $\delta$, against its validity.

For a kernel the obstruction holds even without sampling. For any kernel satisfying (A1) and (A2), let $Q_\nu$ have
density $1+\varsigma(\nu\,\cdot)$ on $[0,1]$, for a measurable $1$-periodic mean-zero $\varsigma$ with $|\varsigma|=1$ and
an integer $\nu\ge1$. Then $\TV(P_{\mathrm u},Q_\nu)=\tfrac12$ for every $\nu$, while $\mu_{Q_\nu}-\mu_{P_{\mathrm u}}=\int_0^1\phi(x)\,
\varsigma(\nu x)\,dx\to0$ in $\Hh$ by the Riemann--Lebesgue lemma for Bochner integrals, since the bounded, strongly
measurable $\phi$ lies in $L^1([0,1];\Hh)$, where step functions are dense, and $\int_a^b\varsigma(\nu x)\,dx\to0$. As
$\mu_P\in\mathcal C_\alpha(P)$ and $\mu_Q\in\mathcal C_\alpha(Q)$, $D_\alpha(P,Q)\le\MMD(P,Q)$ for every $\alpha$, so the
trimmed path is uniformly small while the total variation stays at $\tfrac12$.
